\documentclass[11pt]{article}

\usepackage[margin=1.05in]{geometry}
\usepackage{amsmath,amssymb,mathtools}
\usepackage{booktabs}
\usepackage{xcolor}
\usepackage[colorlinks=true,linkcolor=blue!55!black,citecolor=blue!55!black,urlcolor=blue!55!black]{hyperref}

\newcommand{\F}{\mathbb F_{2^{128}}}
\newcommand{\cube}[1]{\{0,1\}^{#1}}
\newcommand{\eq}{\operatorname{eq}}
\newcommand{\mle}[1]{\widetilde{#1}}
\newcommand{\bits}{\operatorname{bits}}
\emergencystretch=1.5em

\title{\textbf{Jagged polynomial packing}}
\author{}
\date{July 2026}

\begin{document}
\maketitle

\begin{abstract}
\noindent
This note describes the Jagged polynomial commitment used in leanVM. Unequal-height trace columns are packed into one Ligerito commitment without a power-of-two gap after every column, and a width-four branching program evaluates the resulting opening weights. Grand-product leaves retain their simpler aligned power-of-two blocks: the product GKR is accelerated by contracting two tree levels per layer and by keeping its identity-padding suffix implicit. The first construction follows the Basic Jagged method~\cite{Jagged26}; the grand-product optimizations are independent of it.
\end{abstract}

\section{The interval-transport polynomial}

All bit vectors are little-endian. For $x,u\in\F^m$, define
\[
\eq(x,u):=\prod_{i=0}^{m-1}(1+x_i+u_i).
\]
On the Boolean cube this is the equality indicator, so the multilinear extension of $P:\cube m\to\F$ is $\mle P(x)=\sum_{u\in\cube m}P(u)\eq(x,u)$.

Suppose a logical table starts at dense offset $t$ and has $h$ entries. Define
\[
J_{t,h}(r,x):=\sum_{u=0}^{h-1}\eq(r,u)\eq(x,t+u).
\]
At Boolean $x$, this is zero outside $[t,t+h)$, while $J_{t,h}(r,t+u)=\eq(r,u)$. Hence, if $q[t+u]=P[u]$ for $0\leq u<h$, then
\[
\sum_{x\in\cube m}q[x]J_{t,h}(r,x)=\sum_{u=0}^{h-1}P[u]\eq(r,u).
\]
Thus $J$ transports the ordinary equality kernel on a logical table to the table's physical interval inside a larger commitment.

\paragraph{Width-four evaluation.} The verifier evaluates $J$ in $O(m)$ field operations. Scan bits from least to most significant and maintain four weighted states $(\kappa,\chi)\in\cube2$. Here $\kappa$ is the carry in $x=t+u$, while $\chi$ says that the processed prefix of $u$, read with the latest bit as most significant, is strictly smaller than the corresponding prefix of $h$. For Boolean bits $a=u_i$, $d=x_i$, $s=t_i$, and $\ell=h_i$, a valid transition satisfies
\[
d=a+s+\kappa\pmod 2,\qquad \kappa'=\left\lfloor\frac{a+s+\kappa}{2}\right\rfloor,\qquad \chi'=\begin{cases}\chi,&a=\ell,\\ \ell,&a\neq\ell.\end{cases}
\]
The choice $a\in\{0,1\}$ contributes weight $(1+r_i,r_i)$, and the forced output bit $d$ contributes $(1+x_i,x_i)$. Sum the weighted transitions reaching each new state. Start at $(0,0)$, process one extra zero bit above the committed cube, and return state $(0,1)$. The extra bit accepts the exact endpoint $t+h=2^m$ but rejects overflow. This is the Basic Jagged read-once branching program~\cite{Jagged26}.

\section{Jagged trace commitment}

\paragraph{Shift away public padding.} A logical column $P_i$ has $2^{k_i}$ entries: its first $h_i$ entries are witness values and the rest equal a public padding value $p_i$. Over characteristic two define
\[
Q_i[z]=\begin{cases}P_i[z]+p_i,&z<h_i,\\0,&h_i\leq z<2^{k_i}.
\end{cases}
\]
Then $P_i=Q_i+p_i$ on the entire logical cube, so $\mle P_i(r)=\mle Q_i(r)+p_i$. Only the nonzero prefix of $Q_i$ needs committed space, and the verifier never needs a correction for a truncated padding prefix.

\paragraph{Physical layout.} Compatible shifted prefixes of equal height are grouped into a row-major block of width $B=2^b$. If the block starts at $t$, its slot $c$ is stored as
\[
q[t+c+Bz]=Q_c[z],\qquad 0\leq c<B,\quad 0\leq z<h.
\]
The block occupies exactly $Bh$ cells. Placing wider blocks first aligns every start to its width without inserting gaps. After the final block, $q$ is padded once with zero to length $2^M$. Flock's packed witness $q_{\rm pkd}$ remains one unshifted power-of-two block anchored at zero; the Jagged layout is used for the ordinary columns.

\paragraph{Opening and batching.} To open slot $c$ at row point $r$, form $r^\star=(\bits_b(c),r,0,\ldots,0)$ and set $L=Bh$. Then
\[
\sum_xq[x]J_{t,L}(r^\star,x)=\mle Q_c(r),
\]
after which the verifier adds $p_c$. To batch all $B$ slots with coefficients $1,\gamma,\ldots,\gamma^{B-1}$, replace the selector-coordinate equality weights $(1+r_i,r_i)$ inside the branching program by $(1,\gamma^{2^i})$. A slot $c=\sum_i2^ic_i$ then receives $\prod_i\gamma^{2^ic_i}=\gamma^c$. One branching-program evaluation therefore produces the whole geometric batch. The prover materializes the corresponding dense weight vector and uses the unchanged Ligerito opening protocol.

\paragraph{Recursive terminal.} Suppose Ligerito folds the first $f$ committed coordinates and leaves a final message $m:\cube\ell\to\F$, where $f+\ell=M$. Recursion runs the four-state automaton through the folded coordinates and contracts the remaining $2^\ell$ message entries through the residual transition matrices. When every residual logical-row coordinate is fixed to zero, only two adjacent message entries can survive; the implementation selects those entries directly instead of traversing the full residual tree.

\section{Grand-product GKR}

The grand product deliberately does not use the tight Jagged layout above. A bus block $b$ has $2^{\kappa_b}$ rows, including its public default-padding rows. Blocks are sorted by decreasing $\kappa_b$ and placed consecutively. A block of size $2^k$ then starts at a multiple of $2^k$, because every preceding block size is also a multiple of $2^k$. If the total length lies between $2^{\mu-1}$ and $2^\mu$, the single remaining suffix is filled with the product identity $1$.

For block $b$, row $u$, tuple coordinates $c_{b,j}(u)$, and transcript challenges $\alpha,\gamma$, define the leaf
\[
L_b(u)=\gamma+\sum_j\alpha^jc_{b,j}(u).
\]
If its aligned offset is $t_b$, write $s_b=t_b/2^{\kappa_b}$. At a GKR terminal point $\zeta\in\F^\mu$, split $\zeta=(\zeta_{<\kappa_b},\zeta_{\geq\kappa_b})$. The leaf MLE is reconstructed as
\[
\mle L(\zeta)=1+\sum_b\eq\!\left(\zeta_{\geq\kappa_b},\bits(s_b)\right)\left(1+\gamma+\sum_j\alpha^j\mle{c_{b,j}}(\zeta_{<\kappa_b})\right).
\]
The leading $1$ and the inner $1$ account for the outer identity padding in characteristic two. Constants and public columns are evaluated locally; each committed-column term becomes an ordinary PCS opening claim at the prefix $\zeta_{<\kappa_b}$. Repeated claims on the same column and point are sent only once.

Push and pull contain paired blocks and therefore have the same depth. Their roots include different numbers of public default-padding fingerprints, so the verifier removes the known surpluses before comparing them. The count lane is the product of read counts and may have smaller depth; its root must be nonzero.

\subsection{One radix-four proof for three products}

Push, pull, and count are combined with weights $1,\lambda,\lambda^2$ inside one GKR. After every layer a fresh $\lambda$ is sampled, so the individual terminal values cannot be changed while preserving only their old linear combination. Both native and recursive verifiers obtain all three leaf claims at one shared point $\zeta$.

Two binary product-tree levels are contracted at once:
\[
V_i(x)=\prod_{a,b\in\{0,1\}}V_{i-2}(a,b,x).
\]
The equality-trick sumcheck for this relation has degree four. In one round let $q(X)=\sum_{j=0}^4c_jX^j$, let $s$ be the corresponding coordinate of the current claim point, and let $T$ be the incoming claim. Then
\[
T=(1+s)q(0)+sq(1)=c_0+s\bigl(q(0)+q(1)\bigr).
\]
The prover sends only
\[
d=q(0)+q(1),\qquad c_2,\qquad c_3,\qquad c_4.
\]
The verifier recovers $c_0=T+sd$ and $c_1=d+c_2+c_3+c_4$, then evaluates $q$ at the new challenge by Horner's rule. These four elements are exactly the independent data of a degree-four round once the incoming sum is fixed. At the end of the layer the prover sends the four child evaluations for each lane, the verifier checks their products in the current $1,\lambda,\lambda^2$ combination, and two fresh challenges bind the child coordinates. An odd-depth tree begins with its single root-most binary layer.

The prover stores the four child values of each row interleaved. Folding can therefore consume a tree level in place; no full transposition into four separate arrays is needed. The quartic coefficient kernel batches field products at the widest available hardware width: four lanes with AVX-512, two pairs with AVX2/VPCLMUL, two pairs with NEON/PMULL, and a scalar fallback. Reduced and deferred-reduction products have the same architecture coverage.

\subsection{Implicit identity suffix}

Let the count tree have depth $\mu_c\leq\mu$. Logically it is extended to depth $\mu$ by appending $2^\mu-2^{\mu_c}$ leaves equal to $1$, but the prover never allocates or scans them. This optimization is generic: it depends only on a power-of-two prefix followed by identities, not on any particular trace shape.

To see why it works, consider a radix-four sumcheck round after the explicit prefix has shrunk to one row. Every wholly implicit row contains four ones. Its gate polynomial is identically one, so its compact message $(d,c_2,c_3,c_4)$ is $(0,0,0,0)$: $q(0)=q(1)$ in characteristic two and there are no higher coefficients. Only the boundary pair, consisting of the current explicit row and one all-one row, contributes. After sampling $r$, each of its four entries $a$ becomes
\[
\operatorname{interp}(a,1;r)=a+r(a+1).
\]
The same constant-size boundary step repeats for every omitted coordinate. Thus the sparse prover emits exactly the transcript of the dense identity-padded tree, while its work remains proportional to the actual count tree.

\section{Soundness and measured cost}

The Jagged opening is an ordinary linear opening of the committed vector; the branching program only evaluates its public weight. The grand-product GKR keeps the verifier transcript of the dense protocol. Omitting identity rows changes neither a root nor a round polynomial, and the dense-MLE equivalence is tested across even and odd tree depths and several depth gaps. Radix four changes a layer's degree from two to four while halving its number of product levels; the verifier accounts for the degree-four messages explicitly. Bus challenges are sampled after the witness commitment and protected by the existing grinding bound.

On the reference machine with $\texttt{RAYON\_NUM\_THREADS=11}$, $\texttt{recursion --n 2}$ uses $1{,}827{,}834$ guest cycles, proves in $4.72$ seconds, and produces a $593.547$ KiB proof. The corresponding \texttt{jagged-pcs} baseline used $1{,}865{,}318$ guest cycles and about $5.5$ seconds. The command $\texttt{xmss --n-signatures 890}$ proves in $4.52$ seconds with a $590.031$ KiB proof, compared with roughly $4.7$ seconds and $595.422$ KiB on the same baseline. Timings are noisy; the cycle count and protocol sizes are deterministic for these inputs.

Tests compare every SIMD batch primitive with scalar multiplication, compare compact quartic messages with direct polynomial evaluation, compare sparse identity suffixes with dense padding, exercise even and odd GKR depths, and run native and recursive proof round trips.

\begin{thebibliography}{9}

\bibitem{Jagged26}
Tamir Hemo, Kevin Jue, Eugene Rabinovich, Gyumin Roh, and Ron D.\ Rothblum.
\newblock \emph{Jagged Polynomial Commitments (or: How to Stack Multilinears)}.
\newblock Manuscript, 2026.

\end{thebibliography}

\end{document}
